% answers_short.tex —— 2020 年·全国Ⅲ卷·文科 · 简短答案速查
% 由 tools/make_short_answers.py 自 answers.tex 的 \gksolhead/\gkanswer 提取，
% 并与 2020/2020_数学真题_全国3文_答案要点.md 逐题比对后生成，勿手改（改源 MD 后重跑）。

\begin{gknotice}
\textbf{使用说明}：本节只给最终答案便于快速核对，完整推导见解析版。答案为编写者独立推导，\textbf{非官方标准答案}。
\end{gknotice}

\gkpart{选择题}
\begin{center}\begin{tabular}{ccccccccccccc}
\toprule
题号 & 1 & 2 & 3 & 4 & 5 & 6 & 7 & 8 & 9 & 10 & 11 & 12 \\
\midrule
答案 & B & D & C & C & B & A & B & B & C & A & C & D \\
\bottomrule
\end{tabular}\end{center}

\gkpart{填空题}
\begin{itemize}[leftmargin=2.6em]
  \item[\textbf{13.}] $7$。可行域是顶点为 $(0,0)$，$(1,-1)$，$(1,2)$ 的三角形，$z=3x+2y$ 在 $(1,2)$ 处取最大值 $7$。
  \item[\textbf{14.}] $\sqrt3$。$\dfrac ba=\sqrt2$，故 $e=\sqrt{1+\dfrac{b^2}{a^2}}=\sqrt3$。
  \item[\textbf{15.}] $1$。$f'(x)=\dfrac{\mathrm e^x(x+a-1)}{(x+a)^2}$，由 $f'(1)=\dfrac{\mathrm e}4$ 得 $4a=(a+1)^2$，即 $(a-1)^2=0$。
  \item[\textbf{16.}] $\dfrac{\sqrt2}3\pi$。圆锥高 $2\sqrt2$，内切球半径 $\rho=\dfrac{1\times2\sqrt2}{1+3}=\dfrac{\sqrt2}2$，体积 $\dfrac43\pi\rho^3=\dfrac{\sqrt2}3\pi$。
\end{itemize}

\gkpart{解答题}
\begin{itemize}[leftmargin=2.6em]
  \item[\textbf{17.}] (1) $a_n=3^{n-1}$。设公比为 $q$，由 $a_1(1+q)=4$、$a_1(q^2-1)=8$ 得 $q=3$，$a_1=1$；(2) $m=6$。$\log_3a_n=n-1$，$S_n=\dfrac{n(n-1)}2$，由 $S_m+S_{m+1}=S_{m+3}$ 得 $m^2-5m-6=0$，舍去 $m=-1$，得 $m=6$
  \item[\textbf{18.}] (1) $0.43$，$0.27$，$0.21$，$0.09$。四个等级的天数依次为 $43$，$27$，$21$，$9$，各除以 $100$；(2) $350$。$\dfrac{20\times100+35\times300+45\times500}{100}=350$；(3) $K^2\approx5.82>3.841$，有 95\% 把握认为有关（列联表见解析版）
  \item[\textbf{19.}] (1) 证明。向量法得 $\overrightarrow{EF}\cdot\overrightarrow{AC}=AB^2-BC^2=0$，故 $EF\perp AC$；(2) 证明。由 $\overrightarrow{AC_1}=\overrightarrow{AE}+\overrightarrow{AF}$ 得 $C_1$ 在平面 $AEF$ 内
  \item[\textbf{20.}] (1) $k\leqslant0$ 时在 $\mathbf R$ 上单调递增；$k>0$ 时在 $\left(-\infty,-\sqrt{\dfrac k3}\right)$ 与 $\left(\sqrt{\dfrac k3},+\infty\right)$ 上单调递增，在 $\left(-\sqrt{\dfrac k3},\sqrt{\dfrac k3}\right)$ 上单调递减；(2) $\left(0,\dfrac4{27}\right)$。极大值 $k^2+\dfrac23k\sqrt{\dfrac k3}>0$ 恒成立，极小值 $k^2-\dfrac23k\sqrt{\dfrac k3}<0$ 等价于 $k<\dfrac4{27}$
  \item[\textbf{21.}] (1) $\dfrac{x^2}{25}+\dfrac{16y^2}{25}=1$。由 $e^2=\dfrac{25-m^2}{25}=\dfrac{15}{16}$ 得 $m^2=\dfrac{25}{16}$；(2) $\dfrac52$。设 $Q(6,q)$，由垂直且等长得 $P(5-q,1)$ 或 $P(5+q,-1)$，代回椭圆得 $q=\pm2$ 或 $q=\pm8$，四种情形的面积同为 $\dfrac52$
\end{itemize}

\gkpart{选做题}
\begin{itemize}[leftmargin=2.6em]
  \item[\textbf{22.}] (1) $4\sqrt{10}$。令 $y=0$ 得 $t=2$、$A(-4,0)$，令 $x=0$ 得 $t=-2$、$B(0,12)$，故 $|AB|=\sqrt{16+144}=4\sqrt{10}$；(2) $3\rho\cos\theta-\rho\sin\theta+12=0$（即 $\rho(\sin\theta-3\cos\theta)=12$）。直线 $AB$ 的直角坐标方程为 $3x-y+12=0$，代入 $x=\rho\cos\theta$，$y=\rho\sin\theta$ 即得
  \item[\textbf{23.}] (1) 证明。由 $(a+b+c)^2=0$ 得 $ab+bc+ca=-\dfrac{a^2+b^2+c^2}2$，而 $abc=1$ 保证 $a^2+b^2+c^2>0$，故 $ab+bc+ca<0$；(2) 证明。最大值 $a>0$，$b$，$c$ 为 $t^2+at+\dfrac1a=0$ 的两实根，判别式 $a^2-\dfrac4a\geqslant0$，故 $a^3\geqslant4$，即 $a\geqslant\sqrt[3]{4}$
\end{itemize}
